% Ссылка на приложение обязательна в тексте (п. 4.11.4)
% Пример ссылки: (приложение~Б)
\chapter{КРИТЕРИИ ОЦЕНИВАНИЯ ПРОЕКТА}
\label{app:categories}

В приложении перечислены категории оценивания, максимальные баллы и
конкретные критерии, за выполнение которых выставляется полный балл по
соответствующему пункту.

\section{Обработка и подготовка данных}
\label{sec:crit:data}

Максимум 18 баллов (самостоятельный парсинг данных даёт дополнительно до
4 баллов, см. пункт~\ref{crit:data:parsing}).

\begin{enumerate}
  \item Поиск и источник данных: как данные найдены, почему выбран именно
  этот источник.
  \label{crit:data:source}

  \item Описание датасета: объём (число строк и столбцов); если данных
  мало - обоснование достаточности.
  \label{crit:data:dataset}

  \item Полная очистка данных: пропуски, дубликаты, выбросы, корректность
  типов.
  \label{crit:data:cleaning}

  \item Работа с признаками: исходное количество признаков, новые признаки,
  feature engineering.
  \label{crit:data:features}

  \item Визуализации и графики, отражающие зависимости в данных.
  \label{crit:data:viz}

  \item Корректное разбиение на train/validation/test и пояснение, как
  избегали утечки данных (data leakage), если она могла возникнуть.
  \label{crit:data:split}

  \item Выбор метрик качества и обоснование выбора; при нескольких метриках
  - указание, какой метрике отдаётся предпочтение.
  \label{crit:data:metrics}

  \item Самостоятельный парсинг данных (дополнительно до 4 баллов).
  \label{crit:data:parsing}
\end{enumerate}

\section{Моделирование и эксперименты}
\label{sec:crit:modeling}

Максимум 25 баллов.

\begin{enumerate}
  \item Baseline: простая модель <<из коробки>> (линейная или логистическая
  регрессия, k ближайших соседей и аналоги) без feature engineering.
  \label{crit:model:baseline}

  \item Не менее 4-5 различных моделей и ансамблей (случайный лес,
  градиентный бустинг XGBoost/LightGBM и др.).
  \label{crit:model:variety}

  \item Эксперименты: таблица экспериментов с разными моделями, методами и
  перебором гиперпараметров.
  \label{crit:model:experiments}

  \item Уменьшение размерности: при большом числе признаков - эксперименты
  с методами снижения размерности и визуализация результатов.
  \label{crit:model:dimred}

  \item Обоснование выбора итоговой (финальной) модели.
  \label{crit:model:final}
\end{enumerate}

\section{Качество кода и воспроизводимость}
\label{sec:crit:code}

Максимум 12 баллов.

\begin{enumerate}
  \item Чистая структура проекта и понятные имена файлов, модулей и
  сущностей.
  \label{crit:code:structure}

  \item Настроены линтеры (например, ruff, flake8).
  \label{crit:code:linters}

  \item Фиксированный seed (зерно генератора случайных чисел) в
  экспериментах.
  \label{crit:code:seed}

  \item Файл requirements.txt или pyproject.toml с указанием версий зависимостей.
  \label{crit:code:deps}

  \item Контейнеризация: Docker и при необходимости docker-compose.
  \label{crit:code:docker}

  \item В README описана структура проекта и порядок воспроизведения
  результатов.
  \label{crit:code:readme}
\end{enumerate}
